\section{Necessity of the divergence condition}\label{lv:sec:nec}

The sufficiency proof of Theorem~\ref{lv:thm:single} used nothing about $\beta$; the necessity proof divides by $\beta$ on the line, so it needs a hypothesis. Write $K_c(t)=\Gamma(c+it)\beta(c+it)$, which is half the Mellin transform of $\sech$ on the line $\Re s=c$, and let $c>0$ (the spaces $\mathcal W_c$ contain the dilates only for $c>0$).
\begin{quote}
$(\mathrm N_c)$\quad $1/\beta(c+it)\ll(1+|t|)^{A}$ for some $A\ge0$.
\end{quote}

\begin{lemma}\label{lv:lem:Nc}
Let $c>0$. Then $(\mathrm N_c)$ holds unconditionally for $c\ge1$. If $\beta$ has no zeros in $\Re s>\theta$, it holds for every $c>\theta$ and, by the functional equation, for every $c\in(0,1-\theta)$. In particular it holds for all $c>0$ outside $[\frac18,\frac78]$ given a zero free half plane $\Re s>\frac78$ for $L(s,\chi_4)$ (conditional on the quasi Riemann hypothesis of Section~\ref{lv:sec:qrh}), and for all $c>0$, $c\ne\frac12$, under GRH for $\beta$.
\end{lemma}

\begin{proof}
For $c>1$, $|1/\beta(c+it)|\le\zeta(c)$; for $c=1$, $1/\beta(1+it)\ll\log(2+|t|)$ as for $\zeta$. Given a zero free half plane $\Re s>\theta$, the Borel and Carath\'eodory argument of Titchmarsh \cite[Theorem 14.2]{Titchmarsh}, which uses only the zero free region, gives $\log\beta(\sigma+it)\ll(\log|t|)^{B}$ with $B\le1$ for $\sigma>\theta$, hence $(\mathrm N_c)$. For $c<1-\theta$ put $s=1-c-it$, so that $\Re s>\theta$ and the functional equation $\beta(1-s)=(2/\pi)^s\sin(\pi s/2)\Gamma(s)\beta(s)$ gives $\beta(c+it)=(2/\pi)^s\sin(\pi s/2)\Gamma(s)\beta(s)$. Since $|\sin(\pi s/2)\Gamma(s)|\asymp|t|^{1/2-c}$ for large $|t|$ on the line $\Re s=1-c\in(\theta,1)$, which is not an even integer, $1/\beta(c+it)$ is bounded by $|t|^{c-1/2}$ times a polynomial bound for $1/\beta(s)$.
\end{proof}
% UNCHECKED: Titchmarsh Theorem 14.2 for L(s,chi_4); the statement is quoted from the source, the theorem number is not claimed.

\begin{theorem}[necessity]\label{lv:thm:nec}
Let $c>0$ and assume $(\mathrm N_c)$. If $\Lambda\subset[1,\infty)$ and $\sum_\Lambda1/\lambda<\infty$, then $\{\sech(\lambda x)\}_{\lambda\in\Lambda}$ is not complete in $\mathcal W_c$. So under $(\mathrm N_c)$ the dilates are complete in $\mathcal W_c$ if and only if $\sum_\Lambda1/\lambda=\infty$.
\end{theorem}

\begin{proof}
Step 1, a Blaschke product. Put $B(w)=\prod_{\lambda\in\Lambda}\frac{1-w/\lambda}{1+w/\lambda}$. Since $\log\frac{1-w/\lambda}{1+w/\lambda}=-2w/\lambda+O(|w|^2/\lambda^2)$ and $\sum1/\lambda<\infty$, the product converges uniformly on compact subsets of $\C\setminus(-\Lambda)$ to an analytic function with zeros exactly at $\Lambda$ and poles at $-\Lambda$. Each factor has modulus at most $1$ on $\Re w\ge0$ and modulus $1$ on the imaginary axis, so $|B|\le1$ on $\Re w\ge0$. On the imaginary axis
\[
\Bigl|\frac{d^j}{dw^j}\log B(iy)\Bigr|=(j-1)!\Bigl|\sum_\lambda\Bigl[\frac1{(\lambda-iy)^j}+\frac{(-1)^{j-1}}{(\lambda+iy)^j}\Bigr]\Bigr|\le2(j-1)!\sum_\lambda\frac1\lambda,
\]
because $|\lambda\pm iy|\ge\lambda\ge1$. Hence every derivative $B^{(j)}$ is bounded on the imaginary axis.

Step 2, the strip. Let $S=\{z:|\Im z|<\pi/2\}$ and $w=e^z$, which maps $S$ onto the right half plane and $\partial S$ onto the imaginary axis. Put $\widetilde\Phi(z)=B(e^z)e^{-z^2}$. It is analytic on $|\Im z|<\pi$ (the poles of $B(e^z)$ are at $z=\log\lambda+(2j+1)i\pi$), vanishes at $z=\log\lambda$, $\lambda\in\Lambda$, and $|\widetilde\Phi(\xi+i\eta)|\le e^{\pi^2/4-\xi^2}$ on $\overline S$. On the boundary lines $z=\xi\pm i\pi/2$, $\frac{d^k}{d\xi^k}[B(e^z)]=\sum_{j=1}^kS(k,j)e^{jz}B^{(j)}(e^z)$ with Stirling numbers $S(k,j)$, which by Step 1 is $O_k(\sum_{j\le k}e^{j\xi})=O_k(e^{k\max(\xi,0)})$, while the derivatives of $e^{-z^2}$ are polynomials times $e^{-z^2}$. Every $\xi$ derivative of $\widetilde\Phi$ on either boundary line is therefore in $L^2(d\xi)$.

Step 3, the inverse transform. Put $H(t)=\frac1{2\pi}\int\widetilde\Phi(\xi)e^{it\xi}d\xi$. Shifting the contour to $\Im z=\eta$, allowed by the Gaussian decay, gives $H(t)e^{t\eta}=\frac1{2\pi}\int\widetilde\Phi(\xi+i\eta)e^{it\xi}d\xi$ for $|\eta|\le\pi/2$, so $H(t)e^{\pm\pi t/2}$ is the inverse Fourier transform of the boundary values. Since these have all derivatives in $L^2$, Plancherel gives $H(t)e^{\pi|t|/2}(1+|t|)^{A'}\in L^2(\R)$ for every $A'\ge0$, in particular $H(t)e^{\pi|t|/2}(1+|t|)^{A'}\in L^1$, so inversion holds on the closed strip: $\int H(t)e^{-itz}dt=\widetilde\Phi(z)$. Hence $\int H(t)\lambda^{-it}dt=\widetilde\Phi(\log\lambda)=0$ for $\lambda\in\Lambda$, and $H\not\equiv0$ because $\widetilde\Phi\not\equiv0$.

Step 4, division by the kernel. Define $G$ on $\Re s=c$ by $\overline{G(c+it)}=H(t)/K_c(t)$. By Stirling, $|1/\Gamma(c+it)|\asymp e^{\pi|t|/2}|t|^{1/2-c}$ for large $|t|$, and by $(\mathrm N_c)$, $|G(c+it)|\le C|H(t)|e^{\pi|t|/2}(1+|t|)^{A+1/2-c}\in L^2$. Let $g\in\mathcal W_c$ have Mellin transform $G$; $g\ne0$. By \eqref{lv:eq:plancherel} and Lemma~\ref{lv:lem:mellin},
\[
\langle\sech(\lambda\,\cdot),g\rangle_c=\frac{\lambda^{-c}}\pi\int\lambda^{-it}K_c(t)\overline{G(c+it)}\,dt=\frac{\lambda^{-c}}\pi\int H(t)\lambda^{-it}dt=0
\]
for every $\lambda\in\Lambda$.
\end{proof}

\begin{corollary}\label{lv:cor:iff}
For $c\ge1$, for $c>\frac78$ given the zero free half plane $\Re s>\frac78$ for $L(s,\chi_4)$ (conditional), and for every $c>0$, $c\ne\frac12$, under GRH for $\beta$, $\{\sech(\lambda x)\}_{\lambda\in\Lambda}$, $\Lambda\subset[1,\infty)$, is complete in $\mathcal W_c$ if and only if $\sum_\Lambda1/\lambda=\infty$.
\end{corollary}

\begin{remark}[the barrier $c=\frac12$]\label{lv:rem:barrier}
At $c=\frac12$ the line carries the zeros $\frac12+i\gamma$ of $\beta$ and $(\mathrm N_{1/2})$ fails. The construction survives if $H$ can be made to vanish at each $\gamma$ to the order of the zero, with $H/\beta$ still square integrable. Vanishing of $H$ at a point $\gamma$ is the linear condition $\int\widetilde\Phi(\xi)e^{i\gamma\xi}d\xi=0$, a moment condition and not a zero of $\widetilde\Phi$, so the Blaschke product does not produce it; finitely many such conditions could be imposed, but this does not handle all the zeros or the square integrability after division. Whether the zeros of $\beta$ on the critical line can rescue completeness for a set with $\sum1/\lambda<\infty$ is open (Problem~\ref{lv:prob:sharp}). Numerically, for $\Lambda=\{n^2\}$ the distance from $e^{-x}$ to the span decreases only slowly at $c=\frac12$, as it does at $c=\frac32$ where Theorem~\ref{lv:thm:nec} applies, and at $N=24$ it is more than an order of magnitude above the distance for the integers (Table~\ref{lv:tab:squares}). This is consistent with a positive limit, and so with the critical line not being special here; it is not decisive, since incompleteness of a span does not put the particular target $e^{-x}$ outside it.
\end{remark}

\begin{remark}
The argument is the half plane Blaschke construction of M\"untz and Sz\'asz, transported to the strip. What the strip adds is the Gaussian factor $e^{-z^2}$, which supplies the Sobolev regularity of the boundary values needed to absorb the polynomial part of $1/K_c$, and what the kernel adds is the division by $\beta$, which is where $(\mathrm N_c)$ enters. The same proof gives necessity for the zeta kernel $\varphi$ of \eqref{lv:eq:phi}, $0<c<1$, under the corresponding hypothesis $1/\zeta(c+it)\ll(1+|t|)^A$.
\end{remark}

\section{Wiener, Kac and Salem: the $L^1$ picture}\label{lv:sec:salem}

The completeness problems of this chapter have an older $L^1$ form, treated in Section~\ref{wr:sec:wiener} and Section~\ref{wr:sec:salem}: Wiener's theorem \cite{Wiener32} in its homogeneous form \cite{Kac65}, and Salem's criterion \cite{Salem53} for the Fermi kernel $1/(e^x+1)$, whose Mellin transform is $\Gamma(s)\eta(s)$ (Proposition~\ref{wr:prop:wiener}, Corollary~\ref{wr:cor:salem}). In dilation form, for $k$ on $(0,\infty)$ with $\int_0^\infty|k(x)|x^{\sigma-1}dx<\infty$, the span of $\{x^{\sigma-1}k(\lambda x):\lambda>0\}$ is dense in $L^1(0,\infty)$ if and only if the Mellin transform of $k$ has no zero on the line $\Re s=\sigma$, and the bounded annihilators that exist when it vanishes at $\sigma+i\gamma$ are the characters $x^{i\gamma}$ (under $x=e^u$ translates become dilates). The beta analogue for the secant kernel is
\begin{equation}\label{lv:eq:salem-beta}
\overline{\operatorname{span}}\{x^{\sigma-1}\sech(\lambda x):\lambda>0\}=L^1(0,\infty)\iff\beta(\sigma+it)\ne0\ \text{ for all }t\in\R,
\end{equation}
for $\sigma>0$. This holds for $\sigma\ge1$, fails for $\sigma=\frac12$ (the zero $\frac12+6.02\ldots i$), and for $\frac12<\sigma<1$ it holds for every such $\sigma$ if and only if GRH holds for $\beta$; the mirror lines $0<\sigma<\frac12$ carry the same zero free statements by the functional equation. At $\sigma=1$ Salem's own kernel is incomplete, because $1-2^{1-s}$ vanishes at $1+2\pi ik/\log2$ for nonzero integers $k$ (the zero at $k=0$ cancels the pole of $\zeta$, leaving $\eta(1)=\log2\ne0$); the secant kernel has no such defect.

Salem used every positive dilation; the Carleman argument removes that.

\begin{theorem}[Salem's criterion on a divergent set]\label{lv:thm:salem}
Let $\sigma>0$, let $k$ be $1/(e^x+1)$ or $\sech x$, and let $\Lambda\subset[1,\infty)$ satisfy $\sum_{\lambda\in\Lambda}1/\lambda=\infty$. Then $\{x^{\sigma-1}k(\lambda x)\}_{\lambda\in\Lambda}$ is complete in $L^1(0,\infty)$ if and only if the Mellin transform of $k$ has no zero on the line $\Re s=\sigma$. In particular, for $\frac12<\sigma<1$, the integer dilates of the Fermi kernel are complete in $L^1$ if and only if $\zeta$ has no zero on $\Re s=\sigma$, and the integer dilates of $\sech$ if and only if $\beta$ has none.
\end{theorem}

\begin{proof}
If the transform vanishes at $\sigma+i\gamma$ then $x^{i\gamma}$ is a bounded annihilator of every dilate. Conversely let $\varphi\in L^\infty$ annihilate $x^{\sigma-1}k(\lambda x)$ for $\lambda\in\Lambda$ and put $\Phi(w)=\int_0^\infty x^{\sigma-1}\varphi(x)k(wx)\,dx$, $\Re w>0$. Both kernels extend analytically to $\Re w>0$, since $e^{wx}+1$ and $\cosh(wx)$ vanish only for $wx\in i\R$. For $w=re^{i\psi}$, $|\psi|<\pi/2$, split the integral at $x=1/r$. On $x\le1/r$ the kernel is bounded by an absolute constant $C_0$, since $|wx|\le1$ keeps $wx$ away from the zeros. On $x>1/r$, $|e^{wx}+1|\ge e^{rx\cos\psi}-1$ and $|\cosh(wx)|\ge\sinh(rx\cos\psi)\ge\frac12(e^{rx\cos\psi}-1)$, so after $y=rx\cos\psi$ and with a further absolute constant,
\begin{multline*}
|\Phi(w)|\le\|\varphi\|_\infty\Bigl[\frac{C_0}\sigma r^{-\sigma}+C_1(r\cos\psi)^{-\sigma}\int_{\cos\psi}^\infty\frac{y^{\sigma-1}dy}{e^{y}-1}\Bigr]\\
\le C_\sigma\|\varphi\|_\infty\,r^{-\sigma}(\cos\psi)^{-\sigma-1}\log\frac{2}{\cos\psi}
\end{multline*}
for $r\ge1$. So $\Phi(w+1)$ is analytic in $\Re w>0$, and since $|w+1|/\Re(w+1)\le|w|/\Re w$ it is bounded by $C(|w|/\Re w)^{\sigma+2}$. If $\Lambda\cap[1,2)$ is infinite, $\Phi$ vanishes on a set with a limit point in $[1,2]$ and $\Phi\equiv0$ by the identity theorem; otherwise $\Phi(w+1)$ vanishes at $w=\lambda-1$, $\lambda\in\Lambda$, $\lambda\ge2$, a set in $[1,\infty)$ with divergent reciprocal sum. Lemma~\ref{lv:lem:uniq} gives $\Phi\equiv0$. In particular $\varphi$ annihilates $x^{\sigma-1}k(\lambda x)$ for every $\lambda>0$, and the dilation form of Wiener's theorem gives $\varphi=0$.
\end{proof}

Three completeness statements for the same dilates are each governed by which point evaluations of the Mellin transform are continuous on the space.

\begin{center}\footnotesize
\begin{tabular}{lll}
\toprule
space & obstruction & annihilators\\
\midrule
$L^1$, Salem & zeros on $\Re s=\sigma$ & $x^{i\gamma}$\\
$\mathcal W_c$ & none & none\\
$\mathcal W_a\cap\mathcal W_b$ & zeros in $a<\Re s<b$ & $f\mapsto\widehat f(\rho)$\\
\bottomrule
\end{tabular}
\end{center}

In $L^1$ the Mellin transform is continuous on the line and the dual contains the characters, so a zero on the line is seen. In $L^2$ with one weight the transform is only defined almost everywhere on the line, the characters are not square integrable, and nothing off the line is visible: Theorem~\ref{lv:thm:single} holds even at $c=\frac12$, where Salem's family is unconditionally incomplete. With two weights the bound \eqref{lv:eq:pointeval} makes evaluation continuous throughout the strip, and the annihilator attached to a zero $\rho$ is the functional $f\mapsto\widehat f(\rho)$. Salem's criterion and Theorem~\ref{lv:thm:two} are the boundary and interior versions of one statement, detecting zeros on a line and in a strip. The quasi Riemann hypothesis is a statement of the second kind: if it covers $L(s,\chi_4)$, it makes the two weight criterion hold on $(\frac78,1)$ (conditional) but, with only this half plane, does not resolve zero freeness on individual lines in $\frac18\le\Re s\le\frac78$.

\begin{remark}
Nyman and Beurling's criterion is a fourth regime: the dilates are restricted to $(0,1)$, which makes the Mellin transforms of the space analytic in a half plane, so every zero in $\Re s>\frac12$ is seen in a single $L^2$ norm, by the bound $|\int_0^1f(x)x^{s-1}dx|\le\|f\|_2/(2\Re s-1)^{1/2}$ for $\Re s>\frac12$. The space $\mathcal W_a\cap\mathcal W_b$ is the half line substitute for that restriction. Together with Section~\ref{lv:sec:two} this answers the first half of Problem~\ref{rz:prob:nymanbeurling}, in the variable $x$ on $\mathcal W_c$; the $L^2$ version in $s$ on the boundary $\Re\sqrt s=\frac12$ is not addressed.
\end{remark}

\section{Explicit approximants in one weight}\label{lv:sec:gram}

\subsection{The Gram system}

Let $d_N(c)$ be the distance in $\mathcal W_c$ from $e^{-x}$ to the span of $\sech(nx)$, $n\le N$. With $G_{nm}=\langle\sech(n\,\cdot),\sech(m\,\cdot)\rangle_c$ and $b_n=\langle e^{-x},\sech(n\,\cdot)\rangle_c$,
\begin{equation}\label{lv:eq:gram}
d_N(c)^2=\|e^{-x}\|_c^2-b^{\mathsf T}G^{-1}b,\qquad\|e^{-x}\|_c^2=\frac{\Gamma(2c)}{2^{2c}},
\end{equation}
and the optimal coefficients are $a=G^{-1}b$ (the matrix $G$ is positive definite, since a nonzero finite combination of distinct dilates cannot vanish). The data have closed forms. Expanding $\sech(nx)=2\sum_m\chi_4(m)e^{-mnx}$,
\[
b_n=2\Gamma(2c)\sum_{m\ge1}\frac{\chi_4(m)}{(1+nm)^{2c}},\qquad G_{nm}=\frac{2(nm)^{-c}}{\pi}\int_{-\infty}^\infty|K_c(t)|^2\cos\bigl(t\log\tfrac mn\bigr)dt,
\]
the first an alternating series convergent for every $c>0$, the second a rapidly convergent integral because $|K_c(t)|^2\asymp|t|^{2c-1}e^{-\pi|t|}|\beta(c+it)|^2$.

\subsection{What the numbers say}

Table~\ref{lv:tab:dN} gives $d_N(c)$, the largest optimal coefficient, and the error of the M\"obius approximant $S_N$ for comparison, computed in multiple precision (two precisions and step sizes agree to $3\times10^{-28}$ relative in every $d_N$ and $\max|a_n|$; $40$ digits are not enough at $N=32$).

\begin{table}[h]\centering\footnotesize
\begin{tabular}{rrccc}
\toprule
$c$ & $N$ & $d_N(c)$ & $\max_n|a_n|$ & $\|S_N-e^{-x}\|_c$\\
\midrule
$0.60$ & 4 & $8.1\times10^{-3}$ & $0.55$ & $4.8\times10^{-2}$\\
 & 8 & $9.7\times10^{-4}$ & $72$ & $7.3\times10^{-3}$\\
 & 16 & $1.6\times10^{-4}$ & $3.3\times10^{6}$ & $7.8\times10^{-2}$\\
 & 32 & $9.5\times10^{-5}$ & $7.3\times10^{17}$ & $6.2\times10^{-2}$\\
\midrule
$0.75$ & 4 & $5.8\times10^{-3}$ & $0.48$ & $4.0\times10^{-2}$\\
 & 8 & $5.7\times10^{-4}$ & $54$ & $5.1\times10^{-3}$\\
 & 16 & $1.0\times10^{-4}$ & $1.2\times10^{6}$ & $4.7\times10^{-2}$\\
 & 32 & $6.5\times10^{-5}$ & $7.2\times10^{16}$ & $3.5\times10^{-2}$\\
\midrule
$1.00$ & 4 & $3.7\times10^{-3}$ & $0.49$ & $3.0\times10^{-2}$\\
 & 8 & $2.8\times10^{-4}$ & $36$ & $2.9\times10^{-3}$\\
 & 16 & $5.5\times10^{-5}$ & $4.0\times10^{5}$ & $2.1\times10^{-2}$\\
 & 32 & $2.8\times10^{-5}$ & $5.4\times10^{17}$ & $1.4\times10^{-2}$\\
\bottomrule
\end{tabular}
\caption{Optimal distance, size of the optimal coefficients and the M\"obius error (computed).}\label{lv:tab:dN}
% Source discrepancy: the source note printed 2.4e15, 1.5e16, 2.7e16 for max|a_n| at N=32, off by a factor of 5 to 300.
\end{table}

The optimal distance is about one order of magnitude below the M\"obius error at $N=4$ and $8$ (ratios $0.10$ to $0.17$), and two to three orders below it at $N=16$ and $32$ (ratios $1.5\times10^{-3}$ to $2.6\times10^{-3}$). The optimal coefficients explode, between $10^{16}$ and $10^{18}$ at $N=32$: the problem is ill conditioned, and beyond $N\approx16$ the optimal approximant is numerically useless in ordinary precision. The M\"obius approximant, whose coefficients are $0,\pm\frac12$, is stable but converges slowly and not monotonically, its error tracking the oscillation of $M_\chi(N)$. This is the same trade off as in the Nyman and Beurling problem, where B\'aez Duarte \cite{BaezDuarte03} observed that the natural M\"obius approximants are far from optimal.

\subsection{Why the coefficients explode}

The explanation is heuristic. On the line $\Re s=c$ the approximation problem is to make $2\beta D_N-1$ small in $L^2(|\Gamma(c+it)|^2dt)$, where $D_N(c+it)=\sum_{n\le N}a_nn^{-c-it}$. The weight $e^{-\pi|t|}$ confines the problem to a segment $|t|\le T$. On a segment the exponentials $e^{-it\log n}$, $n\ge1$, are complete in $L^2(-T,T)$ for every fixed $T$: if $h$ annihilates them, $E(z)=\int_{-T}^Th(t)e^{-itz}dt$ is entire of exponential type at most $T$, so has $O(r)$ zeros in $|z|\le r$, while $E$ vanishes at the $e^r$ points $\log n\le r$ with $n\le e^r$, so $E=0$. Near linear dependence is quantitative: $\|e^{-it\log(n+1)}-e^{-it\log n}\|_{L^2(-T,T)}^2\le\frac{2T^3}3\log(1+\frac1n)^2$, so the finite Gram matrices have no uniform lower spectral bound. This proves ill conditioning of the restricted systems but does not determine the size of the optimal coefficient vector for this target, and the effective cutoff of the segment depends on the coefficient growth through $\log\bigl((1+\|a\|_{N,c})/\varepsilon\bigr)$, $\|a\|_{N,c}=\sum_{n\le N}|a_n|n^{-c}$. The qualitative picture, completeness together with increasingly close frequencies, is what Table~\ref{lv:tab:dN} shows.

\begin{proposition}[what the weight sees]\label{lv:prop:segment}
Let $c>0$, $T\ge1$ and $a_n$ be arbitrary coefficients, with $D_N(s)=\sum_{n\le N}a_nn^{-s}$ and $\|a\|_{N,c}=\sum_{n\le N}|a_n|n^{-c}$. Then $\|e^{-x}-\sum_{n\le N}a_n\sech(nx)\|_c^2=\frac1{2\pi}\int_{\R}|\Gamma(c+it)|^2|1-2\beta(c+it)D_N(c+it)|^2dt$, and
\[
\begin{aligned}\Bigl\|e^{-x}-\sum_{n\le N}a_n\sech(nx)\Bigr\|_c^2&\le\frac1{2\pi}\int_{|t|\le T}|\Gamma(c+it)|^2\bigl|1-2\beta D_N\bigr|^2dt\\&\quad+C_c\,e^{-\pi T}T^{2c}(1+\|a\|_{N,c})^2 .\end{aligned}
\]
Where $\beta\ne0$ the integrand of the first term is $|K_c(t)|^2|2D_N-1/\beta|^2$.
\end{proposition}

\begin{proof}
The identity is \eqref{lv:eq:plancherel} with Lemma~\ref{lv:lem:mellin}; the Mellin transform of $\sech(nx)$ carries the factor $2$, which is why $2D_N-1/\beta$ and not $D_N-1/\beta$ appears (for $a_n=\mu(n)\chi_4(n)/2$ the first term is the error of $S_N$, which tends to $0$; the form with $D_N-1/\beta$ gives the error of an approximant with twice the coefficients). For the tail $|t|>T$ use $|D_N|\le \|a\|_{N,c}$, $|\Gamma(c+it)|^2\ll_c|t|^{2c-1}e^{-\pi|t|}$ and $|\beta(c+it)|\ll_c(1+|t|)^{1/2}$ (functional equation, absolute convergence to the right of $1$ and interpolation), so the tail is at most $C(1+\|a\|_{N,c})^2\int_T^\infty t^{2c}e^{-\pi t}dt\ll_ce^{-\pi T}T^{2c}(1+\|a\|_{N,c})^2$ for $T\ge1$. The first form holds for every $T>0$ if $T^{2c}$ is replaced by $(1+T)^{2c}$ (take all $a_n=0$ and $T\downarrow0$ to see that some such change is needed).
\end{proof}

A computed check at $N=8$, $c=0.75$, $a_n=\mu(n)\chi_4(n)/2$: the $x$ space error squared is $2.5846230428\times10^{-5}$ and so is $\frac1{2\pi}\int|K_c|^2|2D_N-1/\beta|^2dt$, while the form with $D_N-1/\beta$ gives $7.8\times10^{-2}$.

So an upper bound for $d_N(c)$ follows from approximating $1/\beta$ on a growing segment $|t|\le T_N$ by Dirichlet polynomials with coefficients of size $\|a\|_{N,c}=e^{o(T_N)}$ and a vanishing segment residual (a sufficient route, not an equivalence); the obstruction to an explicit rate is that the coefficient growth is not controlled.

\begin{remark}[what is known conditionally]\label{lv:rem:burnol}
If $\beta$ has no zeros in $\Re s>\theta$, then $d_N(c)\le\|S_N-e^{-x}\|_c=O_\varepsilon(N^{\theta-c+\varepsilon})$ for $c\in(\theta,1)$ by Theorem~\ref{lv:thm:strip}, so $d_N(c)$ decays polynomially; given the quasi Riemann hypothesis for $L(s,\chi_4)$ (conditional) this is a polynomial rate for $c>\frac78$. The question of a lower bound is open. In the Nyman and Beurling setting Burnol \cite{Burnol02} proves a zero dependent asymptotic lower bound of order $1/\log N$ for the squared distance $\delta_N^2$, of the form $\liminf\delta_N^2\log N\ge\sum_{\Re\rho=1/2}m_\rho^2/|\rho|^2$; it comes from the continuity of point evaluation at the zeros, which the single weight $\mathcal W_c$ lacks, and it does not transfer to this Hilbert space. So either there is no polynomial lower bound in one weight or it has a different source; this is the refined form of Problem~\ref{lv:prob:rate}.
\end{remark}
% UNCHECKED: Burnol 2002, the exact form of the lower bound (squared multiplicities, liminf, zeros on the line) is quoted from the reviews of this chapter; the paper was not read for this book. The source printed an inequality for every N with unsquared multiplicities, which is not what was proved.

\section{The Meixner family}\label{lv:sec:meixner}

For $t>0$ let $f_t(x)=|\Gamma(\frac t2+\frac{ix}2)|^2$, $x>0$, the symmetric Meixner density of order $t$ up to normalisation (Section~\ref{rz:sec:meixner}); $f_1(x)=\pi\sech(\pi x/2)$ and $f_2(x)=(\pi x/2)/\sinh(\pi x/2)$.

\begin{theorem}[Mellin transform of the Meixner density]\label{lv:thm:meixner}
For $0<\Re s<1$,
\begin{equation}\label{lv:eq:Mt}
M_t(s)=\int_0^\infty x^{s-1}f_t(x)\,dx=\frac{2^s\pi\,\Gamma(t)}{\sin(\pi s/2)}\,Z_t(1-s),
\end{equation}
where $Z_t$ is entire, given by
\begin{align*}
Z_t(w)&=\frac1{\Gamma(w)}\int_0^\infty u^{w-1}(2\cosh(u/2))^{-t}\,du\\
&=\sum_{k\ge0}(-1)^k\binom{t+k-1}{k}\Bigl(k+\frac t2\Bigr)^{-w}\qquad(\Re w>t-1).
\end{align*} For $t=1,2,3$,
\begin{gather*}
Z_1(w)=2^w\beta(w),\qquad Z_2(w)=\eta(w-1)=(1-2^{2-w})\zeta(w-1),\\
Z_3(w)=2^{w-3}\bigl(\beta(w)-\beta(w-2)\bigr).
\end{gather*} For integer $t$, $Z_t$ is a linear combination of $\eta(w-j)$, $j$ odd, when $t$ is even, and of $2^{w-2i}\beta(w-2i)$ when $t$ is odd.
\end{theorem}

\begin{proof}
The beta integral with $r=\allowbreak e^u$ gives the Fourier representation
\[
|\Gamma(a+iy)|^2=\Gamma(2a)\int_{\R}e^{iyu}(2\cosh(u/2))^{-2a}du,
\]
hence with $a=\frac t2$, $y=\frac x2$,
\[
f_t(x)=2\Gamma(t)\int_0^\infty\cos(xu/2)(2\cosh(u/2))^{-t}du.
\]
Insert an Abel regulator $e^{-\epsilon x}$ in the Mellin integral, so that Fubini applies: the inner integral is
\[
\tfrac12\Gamma(s)[(\epsilon-iu/2)^{-s}+(\epsilon+iu/2)^{-s}],
\]
bounded by $C_su^{-\Re s}$ uniformly in $\epsilon>0$ for $0<\Re s<1$, which is integrable against $(2\cosh(u/2))^{-t}$; dominated convergence as $\epsilon\downarrow0$ gives
\begin{align*}
M_t(s)&=2\Gamma(t)\,2^s\Gamma(s)\cos(\pi s/2)\int_0^\infty u^{-s}(2\cosh(u/2))^{-t}du\\
&=2\Gamma(t)2^s\Gamma(s)\cos(\pi s/2)\Gamma(1-s)Z_t(1-s),
\end{align*}
and
\[
\Gamma(s)\Gamma(1-s)\cos(\pi s/2)=\pi/(2\sin(\pi s/2))
\]
gives \eqref{lv:eq:Mt}. At $u=\allowbreak 0$, $(2\cosh(u/2))^{-t}$ has an even analytic expansion $\sum b_ju^{2j}$; subtracting finitely many terms on $(0,1)$ continues $\int_0^\infty u^{w-1}(\cdots)du$ meromorphically with simple poles at $w=\allowbreak -2j$ only, and division by $\Gamma(w)$ cancels them, so $Z_t$ is entire (and vanishes at $-1,-3,\dots$). Absolute termwise integration of $(2\cosh(u/2))^{-t}=\allowbreak \sum_k\binom{-t}{k}e^{-(k+t/2)u}$ gives the series for $\Re w>t$, and analytic continuation gives it for $\Re w>t-1$, where the alternating series converges. For $t=\allowbreak 1$, $(k+\frac12)^{-w}=\allowbreak 2^w(2k+1)^{-w}$. For $t=\allowbreak 2$, $\binom{k+1}{k}(k+1)^{-w}=\allowbreak (k+1)^{1-w}$. For $t=\allowbreak 3$, 
\[
\binom{k+2}2=\tfrac12\bigl((k+\tfrac32)^2-\tfrac14\bigr)\quad\text{and}\quad\sum_k(-1)^k(k+\tfrac32)^{-u}=2^u(1-\beta(u)),
\]
the constant terms cancelling. For integer $t$, $\binom{k+t-1}{t-1}$ is a polynomial in $m=\allowbreak k+\frac t2$, odd of degree $t-1$ for even $t$ and even for odd $t$, vanishing at
\[
m=1,\dots,\tfrac t2-1,\quad\text{respectively}\quad m=\tfrac12,\dots,\tfrac{t-2}2,
\]
so the sums may be started at $m=1$, respectively $m=\frac12$, and become shifted $\eta$ or $\beta$ series.
\end{proof}

The formulas for $t=2,3$ are elementary and were not checked numerically; the statement for general $t$ is proved above and the explicit combinations are stated for integer $t$ only.

\begin{corollary}[one weight completeness throughout the family]\label{lv:cor:meixner}
For every $t>0$ and every $c>0$, $\{f_t(\lambda x)\}_{\lambda\in\Lambda}$ with $\Lambda\subset[1,\infty)$ is complete in $\mathcal W_c$ whenever $\sum_{\lambda\in\Lambda}1/\lambda=\infty$; in particular the integer dilates are complete.
\end{corollary}

\begin{proof}
$f_t$ is bounded at $0$ and $f_t(x)\sim2\pi(x/2)^{t-1}e^{-\pi x/2}$, so $f_t\in\mathcal W_c$ for all $c>0$. For $c\ge1$ the bounds on $M_t$ below refer to its continuation from $0<\Re s<1$, and at even $c$ the pole of $1/\sin(\pi s/2)$ cancels a zero of $Z_t$. On $\Re s=c$, $|1/\sin(\pi s/2)|\asymp e^{-\pi|\tau|/2}$, and $Z_t(1-c-i\tau)$ is polynomially bounded: with $a_w(k)=\frac{\Gamma(k+t)}{\Gamma(t)\Gamma(k+1)}(k+\frac t2)^{-w}$ and $\Delta a(k)=a(k)-a(k+1)$, alternating summation gives $Z_t(w)=\sum_{j<r}2^{-j-1}\Delta^ja_w(0)+2^{-r}\sum_{k\ge0}(-1)^k\Delta^ra_w(k)$, and on a fixed line $|\Delta^ra_w(k)|\le C(1+|\Im w|)^r(1+k)^{t-1-\Re w-r}$ by Stirling and Cauchy's estimate, so for $r>t-\Re w$ the last series converges absolutely with a polynomial bound in $|\Im w|$. The function $Z_t$ is not identically zero ($M_t(c)>0$), so $M_t(c+i\tau)\ne0$ for almost every $\tau$, and $|M_t(c+i\tau)|\le C(1+|\tau|)^Ae^{-\pi|\tau|/2}$. These are the only properties of the kernel used in the proof of Theorem~\ref{lv:thm:single}.
\end{proof}

The hypothesis $\Lambda\subset[1,\infty)$ cannot be dropped: for $\Lambda=\{2^{-n}:n\ge1\}$, whose reciprocal sum diverges, take a smooth $h\ne0$ supported in a short interval avoiding the zeros of $M_t(c+i\tau)$ together with its translate by $P=2\pi/\log2$, put $H(\tau)=h(\tau)-h(\tau-P)$ and $\overline{G}=H/M_t$; then $\int H(\tau)\lambda^{-i\tau}d\tau=0$ for every $\lambda=2^{-n}$, and $G$ is a nonzero orthogonal vector. The threshold of Remark~\ref{lv:rem:threshold} is universal in the Meixner family: $f_t$ is analytic in the open right half plane, with poles on the imaginary axis at $\pm i(t+2k)$, and the rate $\pi/2$ is the half opening of that sector. What changes with $t$ is the arithmetic.

\begin{theorem}[order two is free]\label{lv:thm:t2}
Let $c_n=\mu(n)$ for odd $n$ and $c_n=0$ for even $n$. Then for every $x>0$, with absolute convergence,
\begin{equation}\label{lv:eq:csch}
e^{-x}=\frac12\sum_{n\ge1}c_n\operatorname{csch}(nx),\qquad xe^{-x}=\frac12\sum_{n\ge1}\frac{c_n}n\,(nx)\operatorname{csch}(nx),
\end{equation}
and $\sum_n(c_n/n)n^{-s}=1/\bigl((1-2^{-s-1})\zeta(s+1)\bigr)$ converges absolutely for $\Re s>0$. Consequently the series for $xe^{-x}$ converges absolutely in every $\mathcal W_c$, $c>0$, and $\{(nx)\operatorname{csch}(nx)\}_{n\in\mathbb N}$ is complete in every $\mathcal W_c$ and in every $\mathcal W_a\cap\mathcal W_b$, $0<a<b$, unconditionally.
\end{theorem}

\begin{proof}
(Equivalently $c_n=\sum_{2^j\mid n}\mu(n/2^j)$: for $n=2^km$ with $m$ odd and $k\ge1$ the two nonzero terms $\mu(m)$ and $\mu(2m)=-\mu(m)$ cancel.) Since $\operatorname{csch}(nx)=2\sum_{m\text{ odd}}e^{-mnx}$, absolute convergence at fixed $x>0$ permits grouping by $mn=N$: $\frac12\sum_nc_n\operatorname{csch}(nx)=\sum_{N\text{ odd}}e^{-Nx}\sum_{n\mid N}\mu(n)=e^{-x}$. Multiply by $x$ for the second identity. The Mellin transform of $(nx)\operatorname{csch}(nx)$ is $2n^{-s}\Gamma(s+1)(1-2^{-s-1})\zeta(s+1)$ for $\Re s>0$. Put $h(x)=x\operatorname{csch}x\in\mathcal W_c$; then $\sum|c_n/n|\,\|h(n\cdot)\|_c\le\|h\|_c\sum n^{-c-1}<\infty$, so the series converges absolutely in each $\mathcal W_c$, and simultaneously in $\mathcal W_a\cap\mathcal W_b$, to its pointwise value $xe^{-x}$. The span is invariant under integer dilations, so its closure contains every $nxe^{-nx}$. These are complete: if $g$ is orthogonal to them in $\mathcal W_a\cap\mathcal W_b$ (or $\mathcal W_c$) and $w(x)=x^{2a-1}+x^{2b-1}$ (respectively $x^{2c-1}$), the finite complex measure $d\nu=xe^{-x}\bar g\,w\,dx$ has $\int e^{-(n-1)x}d\nu=0$ for $n\ge1$; pushing forward by $y=e^{-x}$ and using the density of polynomials in $C[0,1]$ gives $\nu=0$, hence $g=0$.
\end{proof}

The function $M_2(s)=2(2/\pi)^s\Gamma(s+1)(1-2^{-s-1})\zeta(s+1)$ has no zeros in $\Re s>0$; the zeros of $Z_2(1-s)=\eta(-s)$ at the positive even integers cancel the apparent poles of $1/\sin(\pi s/2)$, and the nontrivial zeros of $\zeta$ mapped by $s\mapsto-s$ lie to the left of the admissible half plane (the critical line of $\eta(-s)$ is $\Re s=-\frac12$). There is nothing for the two weight criterion to detect, and Theorem~\ref{lv:thm:t2} is the Meixner analogue of the classical range $c\ge1$ for the secant, now covering every weight. For integer $t\ge3$ the zero set of $Z_t(1-s)$ in the strip is that of a combination such as $\beta(1-s)-\beta(-1-s)$ or $\eta(-s)-\eta(-2-s)$; it is not governed by a single $L$ function, and for noninteger $t$ the series $Z_t$ remains the object to study.

\section{The Thorin form of the two weight criterion}\label{lv:sec:thorinform}

Let $\Xi$ denote $\xi$ or $\xi(\cdot,\chi_4)$, with $\xi(s,\chi_4)=(4/\pi)^{(s+1)/2}\Gamma(\frac{s+1}2)\beta(s)$ as in Section~\ref{sec:dbeta:beta}. Both are entire of order one, real and positive on the real axis, satisfy $\Xi(s)=\Xi(1-s)$, and have the Hadamard product $\Xi(s)=\Xi(0)\prod_\rho(1-s/\rho)$ over the nontrivial zeros, with $\rho$ and $\bar\rho$ paired (Davenport \cite{Davenport}); grouping the zeros into reflected quartets makes the product invariant under $s\mapsto1-s$, so the functional equation forces the exponential factor of Hadamard's theorem to be trivial. For $\theta\ge\frac12$ put
\begin{equation}\label{lv:eq:Ftheta}
F_\theta(s)=\frac{\Xi(\theta)}{\Xi(\theta+\sqrt s)},\qquad s\in\cut,
\end{equation}
with the principal square root. A probability law on $[0,\infty)$ is a generalised gamma convolution (GGC) if its Laplace transform is $\exp(-as-\int_0^\infty\log(1+s/u)U(du))$ with $a\ge0$ and a Thorin measure $U$, a positive Radon measure on $(0,\infty)$ with $\int_{(0,1)}|\log u|U(du)<\infty$ and $\int_1^\infty u^{-1}U(du)<\infty$ (local finiteness is part of the definition: the two integrals alone allow infinite mass accumulating at $1$); the Thorin measure is unique \cite{Bondesson}.

For $\Xi=\xi$ and $\theta\ge1$, and more generally when $\theta>\Re\rho$ for every zero, the GGC statement and the Thorin density are Theorem~\ref{thm:poisson} of Chapter~\ref{chap:evaluation} (with the density \eqref{eq:Ualpha} and its Poisson structure, Proposition~\ref{prop:struct}); the strip form is Theorem~\ref{sm:thm:thorin} and Theorem~\ref{sm:thm:strip}, and Chapter~\ref{chap:basepoint} computes the Thorin measure at the basepoint one. New here are the case of $\xi(s,\chi_4)$, the description of one conjugate pair of zeros as the law of $|\gamma+aC|$, and the link with the two weight criterion of Theorem~\ref{lv:thm:two}.

\begin{lemma}[one pair of zeros]\label{lv:lem:pair}
Let $a\ge0$ and $\gamma>0$. For $s\in\cut$,
\[
\log\frac{(a+\sqrt s)^2+\gamma^2}{a^2+\gamma^2}=\int_0^\infty\log\Bigl(1+\frac su\Bigr)Q_{a,\gamma}(du),
\]
where $Q_{0,\gamma}=\delta_{\gamma^2}$ and, for $a>0$, $Q_{a,\gamma}$ is the probability measure on $(0,\infty)$ with density $q_{a,\gamma}(u)=\frac1{2\pi\sqrt u}\bigl[\frac a{a^2+(\sqrt u-\gamma)^2}+\frac a{a^2+(\sqrt u+\gamma)^2}\bigr]$. In the variable $v=\sqrt u$, $Q_{a,\gamma}$ is the law of $|\gamma+aC|$ with $C$ standard Cauchy.
\end{lemma}

\begin{proof}
Both sides vanish at $s=0$, so it suffices to compare $s$ derivatives for $s>0$ and continue analytically (the positivity of the measure and the bounds below justify differentiation under the integral). The left side has derivative $(a+\sqrt s)/(\sqrt s((a+\sqrt s)^2+\gamma^2))$. With $u=v^2$ the right side has derivative $\frac1\pi\int_{\R}\frac1{v^2+s}\frac a{a^2+(v-\gamma)^2}dv$, the Poisson integral of $v\mapsto1/(v^2+s)$ at the point $\gamma+ia$ of the upper half plane. Since $1/(v^2+s)=\Re[1/(\sqrt s(\sqrt s-iv))]$ and $z\mapsto1/(\sqrt s(\sqrt s-iz))$ is bounded and analytic in $\Im z>0$, the Poisson integral equals $\Re[1/(\sqrt s(\sqrt s+a-i\gamma))]$, which is the left derivative. The mass is $1$ because the two Cauchy densities integrate over $v>0$ to the integral of one Cauchy density over $\R$; the case $a=0$ is the limit. The condition $a\ge0$ is indispensable.
\end{proof}

\begin{theorem}[Thorin form]\label{lv:thm:thorin}
Let $\theta\ge\frac12$. The following are equivalent.
\begin{enumerate}
\item[(i)] $\Xi$ has no zeros in $\Re s>\theta$.
\item[(ii)] $F_\theta$ is the Laplace transform of a generalised gamma convolution.
\item[(iii)] $F_\theta$ is analytic in $\cut$.
\end{enumerate}
When they hold, the Thorin measure of $F_\theta$ is
\begin{equation}\label{lv:eq:Utheta}
U_\theta=\sum_{\rho,\ \Im\rho>0}m_\rho\,Q_{\theta-\Re\rho,\ \Im\rho},
\end{equation}
over the distinct zeros in the upper half plane with multiplicity $m_\rho$, and $F_\theta(s)=\exp(-\int_0^\infty\log(1+s/u)U_\theta(du))$ with no linear term. For $\Xi=\xi(\cdot,\chi_4)$, (i) is equivalent to two weight completeness of the secant dilates on $(\theta,1)$ (Theorem~\ref{lv:thm:two}, for $\theta<1$); for $\Xi=\xi$, to two weight completeness of the dilates of $\varphi$ (Theorem~\ref{lv:thm:zeta}).
\end{theorem}

\begin{proof}
(ii) implies (iii) because a GGC Laplace transform is analytic in $\cut$. (iii) implies (i): a zero $\rho$ with $\Re\rho>\theta$ makes $F_\theta$ have a pole at $s=(\rho-\theta)^2$, which is not on $(-\infty,0]$ since $\Re(\rho-\theta)>0$ (and $\Xi(\theta)>0$, the zero is nonreal and the square root map is locally invertible there). (i) implies (ii): by the Hadamard product and the pairing, for $s>0$,
\[
\begin{gathered}\log\frac{\Xi(\theta+\sqrt s)}{\Xi(\theta)}=\sum_{\Im\rho>0}m_\rho\log\frac{(a_\rho+\sqrt s)^2+\gamma_\rho^2}{a_\rho^2+\gamma_\rho^2},\\ a_\rho=\theta-\Re\rho\in[0,\theta],\ \gamma_\rho=\Im\rho,\end{gathered}
\]
with $a_\rho\ge0$ by (i); the summands are nonnegative and $O_{\theta,s}(\gamma_\rho^{-2})$, so the series converges normally. Lemma~\ref{lv:lem:pair} turns each term into $\int\log(1+s/u)Q_{a_\rho,\gamma_\rho}(du)$, and by Tonelli the sum of the positive measures is a Thorin measure: for large $\gamma$, $0\le a\le\theta$ and $0<u<1$, $q_{a,\gamma}(u)\le C_\theta\gamma^{-2}u^{-1/2}$, which makes the small $u$ integral summable (finitely many low zeros cost nothing), and $\int_1^\infty u^{-1}Q_{a,\gamma}(du)$ is the integral of $v^{-2}$ against the folded Cauchy law over $v\ge1$, which is $\ll_\theta\gamma^{-2}$ (on $|v-\gamma|<\gamma/2$ use $v^{-2}\le4\gamma^{-2}$ and mass at most one, elsewhere the density is at most $C_\theta\gamma^{-2}$ and $\int_1^\infty v^{-2}dv<\infty$; for $a=0$ the atom at $\gamma^2$ gives the same). Hence $U_\theta$ is a locally finite Thorin measure, $F_\theta(s)=\exp(-\int\log(1+s/u)U_\theta(du))$ for $s>0$, and both sides continue to $\cut$. The absence of a linear term is the absence of an exponential factor in the Hadamard product. The completeness statements are Theorems~\ref{lv:thm:two} and~\ref{lv:thm:zeta}.
\end{proof}

\begin{corollary}[Thorin representations above the zero free abscissa]\label{lv:cor:uncond}
For every $\theta\ge1$, $\xi(\theta)/\xi(\theta+\sqrt s)$ and $\xi(\theta,\chi_4)/\xi(\theta+\sqrt s,\chi_4)$ are Laplace transforms of GGCs with the absolutely continuous Thorin measures \eqref{lv:eq:Utheta}, since the zeros have $\Re\rho<1$. Conditionally on the zero free half plane $\Re s>\frac78$ of the quasi Riemann hypothesis (for $\beta$, on the theorem covering $L(s,\chi_4)$), the GGC representation holds for every $\theta\ge\frac78$, and the Thorin measure is absolutely continuous for $\theta>\frac78$, and at $\theta=\frac78$ if no zero lies on $\Re s=\frac78$; a zero on that line contributes an atom $m_\rho\delta_{\gamma_\rho^2}$. For $\theta=\frac12$ the statement for $\xi$ is the Riemann hypothesis and the Thorin measure is then $\sum_{\gamma>0}m_\gamma\delta_{\gamma^2}$.
\end{corollary}

The last sentence is the reciprocal Xi representation of the book; Theorem~\ref{lv:thm:thorin} embeds it in a one parameter family and describes the family above the zero free abscissa.

\begin{theorem}[the Cauchy semigroup of Thorin measures; for $\xi$ the Cauchy flow of Proposition~\ref{prop:struct} and {\cite[Chapter 8]{P26}}]\label{lv:thm:semigroup}
Suppose $\Xi$ has no zeros in $\Re s>\theta_0$. For $\theta\ge\theta_0$ let $\widetilde U_\theta=\frac12\sum_\rho m_\rho\,\mathrm{Cauchy}(\Im\rho,\theta-\Re\rho)$ be the symmetric measure on $\R$ (in the variable $v$), summed over the distinct zeros of both half planes, with $\mathrm{Cauchy}(\gamma,0)=\delta_\gamma$. Then $U_\theta$ is the image of $\widetilde U_\theta$ under $v\mapsto v^2$, and for $\theta_0\le\theta<\theta'$,
\[
\widetilde U_{\theta'}=\widetilde U_\theta*\mathrm{Cauchy}(0,\theta'-\theta).
\]
In particular, if all zeros lie on $\Re s=\theta_0$ (which for the completed functions forces $\theta_0=\frac12$), then $\widetilde U_\theta$ is the Poisson smoothing at width $\theta-\theta_0$ of $\frac12\sum_\rho m_\rho\delta_{\Im\rho}$, and $\widetilde U_\theta\to\frac12\sum_\rho m_\rho\delta_{\Im\rho}$ vaguely as $\theta\downarrow\theta_0$.
\end{theorem}

\begin{proof}
Each of $\mathrm{Cauchy}(\gamma,a)$ and $\mathrm{Cauchy}(-\gamma,a)$ is carried by $v\mapsto v^2$ to $Q_{a,\gamma}$ (Lemma~\ref{lv:lem:pair}), so the factor $\frac12$ makes the image equal to \eqref{lv:eq:Utheta}; without it the image of the sum over all zeros is $2U_\theta$. The Cauchy laws form a convolution semigroup in the scale, $\mathrm{Cauchy}(\gamma,a)*\mathrm{Cauchy}(0,b)=\mathrm{Cauchy}(\gamma,a+b)$, and $\theta'-\Re\rho=(\theta-\Re\rho)+(\theta'-\theta)$; Tonelli permits summing the positive measures, which are locally finite by the tail estimate of Theorem~\ref{lv:thm:thorin}. For vague convergence, a compactly supported continuous test function sees each Cauchy term converge to its atom, and the contribution of large $|\gamma|$ is bounded by $C\sum\gamma^{-2}$ uniformly for widths in a bounded interval. The measures have infinite total mass, so weak convergence is not the notion.
\end{proof}

Three remarks on what this gives the programme. First, a checkable statement. For $\theta\ge1$ the Thorin measure of $\xi(\theta)/\xi(\theta+\sqrt s)$ is the Cauchy smoothing at width $\theta-\frac12$ of the zero measure, plus a correction from any zeros off the critical line. For $\theta=1$ and the first $1000$ zeros, taken on the critical line, the sum $\sum_\gamma\int\log(1+s/u)Q_{1/2,\gamma}(du)$ plus a smooth tail from the zero density reproduces $\log\xi(1+\sqrt s)/\xi(1)$ (Table~\ref{lv:tab:thorin}, computed). The tail is a model and not a bound, so the totals are computed, not certified. A Thorin measure obtained by any other computational route must agree with this one; the comparison at level $\theta\ge1$ is a test of the computational machinery with no hypothesis attached.

\begin{table}[h]\centering\footnotesize
\begin{tabular}{rcccc}
\toprule
$s$ & $\log\xi(1+\sqrt s)/\xi(1)$ & zero sum & tail (model) & total\\
\midrule
1 & 0.0461176 & 0.044678 & 0.00144 & 0.046118\\
4 & 0.137917 & 0.133599 & 0.00432 & 0.137917\\
16 & 0.454853 & 0.440461 & 0.0144 & 0.454854\\
\bottomrule
\end{tabular}
\caption{The Thorin check at $\theta=1$ (computed). The first column is the closed forms $\log(\pi/3)$, $\log(3\zeta(3)/\pi)$, $\log(15\zeta(5)/\pi^2)$ evaluated.}\label{lv:tab:thorin}
\end{table}

Second, the right metric. Above the zero free abscissa the Thorin measure is absolutely continuous and the RH measure $\sum\delta_{\gamma^2}$ is atomic, and both have infinite mass, so their total variation distance is infinite and not $2$; the meaningful comparison is vague convergence (Theorem~\ref{lv:thm:semigroup}), or a metric after a weight with finite integral. The idea that a zero to the right of $\theta$ shows up as a negative bump of mass $-1$ at $\gamma$, by putting $a<0$ in $Q_{a,\gamma}$, is false as a characterisation and is not used: formula \eqref{lv:eq:Utheta} was proved only for $a\ge0$. For $a=-h<0$ one has $L_{-h}(s)=-L_h(s)+\log\frac{(s+h^2+\gamma^2)^2-4h^2s}{(h^2+\gamma^2)^2}$, where $L_a(s)=\log\frac{(a+\sqrt s)^2+\gamma^2}{a^2+\gamma^2}$; a negative Cauchy density accounts for $-L_h$ only, and the last term has its zeros at $(h\pm i\gamma)^2$, off the cut, which are precisely the singularities that obstruct the GGC representation. At $\theta=\frac12$ an off critical reflected pair with widths $-h$ and $+h$ has formally cancelling densities, while $F_{1/2}$ still has off cut poles; so positivity of the formal signed density does not characterise RH (compare Remark~\ref{rem:ch2:oneline}, where a positive density with a zero to its right is exhibited). What is proved is Theorem~\ref{lv:thm:thorin}: RH is equivalent to analyticity of $F_{1/2}$ on $\cut$, or to its GGC property, and above a zero free boundary the positive Cauchy mixture is the genuine Thorin measure. No theorem here identifies the first failure of the GGC property with the appearance of a negative part.

Third, the flow. Above the zero free abscissa $\theta\mapsto\widetilde U_\theta$ is the Poisson extension of the already smoothed boundary measure $\widetilde U_{\theta_0}$, a Cauchy semigroup running forward, and it is the Poisson extension of an atomic zero measure on one line only when all real parts agree. The two weight criterion of Theorem~\ref{lv:thm:two} is, through Theorem~\ref{lv:thm:thorin}, the statement that $F_\theta$ is a GGC at every level above $\theta$. Forward Cauchy convolution is a positive operation; the inverse is not automatically one, so running the semigroup downward from a computed $\widetilde U_1$ is an inversion problem, not a theorem (Problem~\ref{lv:prob:down}).

\section{Numerical checks of the sharpness experiment}\label{lv:sec:numerics}

The sharpness experiment compares $\Lambda=\{n^2\}$ with $\Lambda=\{n\}$ at $c=\frac32$ and at $c=\frac12$ (computed with the Gram system of Section~\ref{lv:sec:gram}).

\begin{table}[h]\centering\footnotesize
\begin{tabular}{rcccc}
\toprule
& \multicolumn{2}{c}{$c=3/2$} & \multicolumn{2}{c}{$c=1/2$}\\
$N$ & squares & integers & squares & integers\\
\midrule
4 & $9.9\times10^{-3}$ & $1.9\times10^{-3}$ & $1.2\times10^{-2}$ & $1.0\times10^{-2}$\\
8 & $6.5\times10^{-3}$ & $1.3\times10^{-4}$ & $7.2\times10^{-3}$ & $1.4\times10^{-3}$\\
16 & $5.2\times10^{-3}$ & $1.6\times10^{-5}$ & $4.9\times10^{-3}$ & $2.4\times10^{-4}$\\
24 & $4.8\times10^{-3}$ & $9.3\times10^{-6}$ & $4.2\times10^{-3}$ & $1.2\times10^{-4}$\\
\bottomrule
\end{tabular}
\caption{Distance from $e^{-x}$ to the span of the first $N$ dilates, squares against integers (computed).}\label{lv:tab:squares}
\end{table}

For the squares the distance flattens, consistent with Theorem~\ref{lv:thm:nec} at $c=\frac32$ (the theorem gives a nonzero orthogonal vector, not that $e^{-x}$ lies outside the closed span), and it behaves similarly at $c=\frac12$; for the integers it keeps falling. The ratio of squares to integers is $5$, $49$, $322$, $519$ at $c=\frac32$ and $1.1$, $5.0$, $20$, $35$ at $c=\frac12$, so at $c=\frac12$ it exceeds one order of magnitude from $N=16$ on and does not reach two.

\section{Open problems}\label{lv:sec:open}

\begin{problem}[sharpness at $c=\frac12$]\label{lv:prob:sharp}
Does every $\Lambda\subset[1,\infty)$ with $\sum1/\lambda<\infty$ admit a nonzero annihilator of $\{\sech(\lambda x)\}$ in $\mathcal W_{1/2}$, that is, can the zeros of $\beta$ on the critical line be used to build, from a convergent set, functions that the Blaschke construction cannot see? A separate target question is whether $e^{-x}$ has positive distance from the closed span of $\{\sech(n^2x)\}$ in $\mathcal W_{1/2}$. The numerical evidence for the squares (Table~\ref{lv:tab:squares}) is consistent with a positive limit and proves nothing.
\end{problem}

\begin{problem}[rate and lower bound for $d_N(c)$]\label{lv:prob:rate}
Find an unconditional rate for $d_N(c)$, $\frac12<c<1$, or a lower bound. By Proposition~\ref{lv:prop:segment} a sufficient route to an upper bound is approximating $1/\beta$ on the segment $|t|\le T_N$ by Dirichlet polynomials with coefficients at most $e^{o(T_N)}$.
\end{problem}

\begin{problem}[the strip converse]\label{lv:prob:strip}
Does a zero free strip $a<\Re s<b$ alone give completeness of $\{\sech(nx)\}$ in $\mathcal W_a\cap\mathcal W_b$? Proposition~\ref{lv:prop:nec} is the converse direction, and Theorem~\ref{lv:thm:two} needs the whole half plane.
\end{problem}

\begin{problem}[Meixner order $t\ge3$]
Describe the zero sets of $Z_t(1-s)$ in the strip for integer $t\ge3$, beginning with $\beta(1-s)-\beta(-1-s)$. Are there infinitely many zeros in $0<\Re s<1$, and is there a two weight criterion with an arithmetic meaning?
\end{problem}

\begin{problem}[below the zero free abscissa]\label{lv:prob:down}
Is there a direct characterisation, not passing through the zeros, of which positive measures $\widetilde U_1$ arise at level $\theta=1$ from an $L$ function? Given $\widetilde U_1$, does the level one Thorin measure together with an explicit analytic reconstruction determine the first abscissa at which $F_\theta$ develops an off cut pole or loses the GGC property? The level one function $\Xi(1+\sqrt s)$ is available from the absolutely convergent Euler product, but recovering $\widetilde U_1$ from it is an ill posed inversion.
\end{problem}
